\section{Casos de Uso (IEEE 830)}

Esta sección provee la especificación detallada de los Casos de Uso, definiendo el comportamiento del sistema bajo diversas condiciones de interacción. La Figura \ref{fig:casos_uso} modela a los actores involucrados.

\begin{figure}[ht]
\centering
\begin{tikzpicture}[
    scale=0.75, transform shape,
    actor/.style={draw, circle, fill=white, minimum size=0.5cm},
    usecase/.style={draw, ellipse, fill=blue!10, text width=2.5cm, align=center, minimum height=1cm},
    system/.style={rectangle, draw, thick, minimum width=6cm, minimum height=5cm}
]

% Sistema
\node[system] (sys) at (3, -1.5) {};
\node[above=0.1cm of sys] {\textbf{Software IA}};

% Actores
\node (student) at (-1, 0) {\textbf{Usuario}};
\node[below=0cm of student] {Estudiante};

\node (api) at (7, 0) {\textbf{API Groq}};

% Casos de Uso
\node[usecase] (uc1) at (3, 0) {Formular Problema};
\node[usecase] (uc2) at (3, -1.5) {Generar Objetivos};
\node[usecase] (uc3) at (3, -3) {Buscar Antecedentes};

% Conexiones
\draw (student) -- (uc1);
\draw (student) -- (uc2);
\draw (student) -- (uc3);

\draw[dashed, ->] (uc1) -- node[above] {\footnotesize \textit{<<invoke>>}} (api);
\draw[dashed, ->] (uc2) -- node[above] {\footnotesize \textit{<<invoke>>}} (api);
\draw[dashed, ->] (uc3) -- node[above] {\footnotesize \textit{<<invoke>>}} (api);

\end{tikzpicture}
\caption{Diagrama de Casos de Uso con Extensión de Actor Remoto}
\label{fig:casos_uso}
\end{figure}

\subsection{Especificación: Formular Objetivos de Investigación (CU-02)}

A continuación, se detalla formalmente el escenario nominal del Caso de Uso \textbf{Generar Objetivos} (modelado como CU-02).

\begin{itemize}
    \item \textbf{Actor Principal:} Estudiante / Tesista.
    \item \textbf{Actor Secundario:} API Groq (Sistema Experto LLM).
    \item \textbf{Precondiciones:} El usuario debe tener conexión a internet activa y debe encontrarse en la vista "Objetivos".
    \item \textbf{Flujo Principal (Nominal):}
    \begin{enumerate}
        \item El Estudiante introduce el \textit{Tema de Investigación} (Ej: "IA en Educación").
        \item El Estudiante ingresa el \textit{Propósito del Estudio} (Ej: "Evaluar el impacto de...").
        \item El Estudiante hace clic en el botón "Generar Objetivos".
        \item El Sistema (Frontend) verifica que los campos obligatorios no estén vacíos.
        \item El Sistema invoca al Actor Secundario (API Groq) enviando el payload.
        \item La API Groq procesa los datos y devuelve el texto con verbos taxonómicos.
        \item El Sistema renderiza los resultados en el componente de Markdown.
    \end{enumerate}
    \item \textbf{Flujos Alternativos (Errores y Excepciones):}
    \begin{itemize}
        \item \textbf{4a. Formulario Incompleto:} Si los campos están vacíos, el sistema despliega un \textit{Toast} (notificación visual) en rojo indicando el error. El flujo retorna al paso 1 sin comunicarse con la API.
        \item \textbf{5a. Falla de Red / Rate Limit:} Si el servidor de Groq no responde (HTTP 500) o está saturado (HTTP 429), el bloque \texttt{catch} interceptará la promesa fallida y mostrará el mensaje "Error de Comunicación: Por favor intente en unos minutos". El estado del botón volverá a estar habilitado.
    \end{itemize}
    \item \textbf{Postcondiciones:} El usuario visualiza una propuesta de 1 Objetivo General y 3 Objetivos Específicos formateados y estructurados.
\end{itemize}

\subsection{Especificación: Formular Problema de Investigación (CU-01)}
\begin{itemize}
    \item \textbf{Actor Principal:} Estudiante / Tesista.
    \item \textbf{Actor Secundario:} API Groq (Sistema Experto LLM).
    \item \textbf{Precondiciones:} Conexión a internet activa; vista "Problema" seleccionada.
    \item \textbf{Flujo Principal:}
    \begin{enumerate}
        \item El Estudiante ingresa el \textit{Tema} y el contexto sintomático del problema.
        \item Selecciona la opción "Generar Planteamiento".
        \item El Sistema valida los datos y bloquea el botón para evitar múltiples envíos.
        \item El Sistema envía el \textit{payload} al modelo LLM.
        \item La API Groq responde con la estructura del problema (Síntomas, Causas, Pronóstico y Control).
        \item El Sistema muestra el resultado en pantalla en formato Markdown.
    \end{enumerate}
    \item \textbf{Flujos Alternativos:} 
    \begin{itemize}
        \item \textbf{Falta de contexto:} Si el LLM considera que la entrada del usuario es demasiado vaga o incoherente, devolverá un texto sugiriendo al usuario reformular su idea inicial.
    \end{itemize}
\end{itemize}

\subsection{Especificación: Buscar Antecedentes (CU-03)}
\begin{itemize}
    \item \textbf{Actor Principal:} Estudiante / Tesista.
    \item \textbf{Actor Secundario:} API Groq y Base de datos bibliográfica.
    \item \textbf{Precondiciones:} El usuario debe tener un tema definido.
    \item \textbf{Flujo Principal:}
    \begin{enumerate}
        \item El Estudiante introduce las palabras clave de su investigación.
        \item Selecciona "Generar Antecedentes".
        \item El Sistema procesa la solicitud mediante \texttt{groqService}.
        \item La API Groq devuelve un resumen estructurado del estado del arte, sugiriendo posibles líneas teóricas históricas (máximo de 5 años de antigüedad según la directriz del *System Prompt*).
        \item El Sistema renderiza los párrafos y viñetas correspondientes.
    \end{enumerate}
\end{itemize}
